\section{Exchange cycles, pins, and zero-free traces}\label{sec:exchange}

The exchange representation gives two finite-dimensional tools for the
proof.  Pinning and exposing cycles turns a return probability into the
first moment of a Markov hitting kernel.  Stability of the same exchange
operators also gives zero-free traces after insertion of suitable
single-slot matrices.  We develop both tools here, keeping track of the
conditioning in the first and the trace contraction in the second.

\subsection{Symmetric slots and permutation layers}

Put $\ell=2S$ and let
$\mathcal I=\Lambda_L\times\{1,\ldots,\ell\}$ be the set of
$N=\ell V$ slots.  On $(\C^2)^{\otimes\mathcal I}$ write
$s_i^a$ for the Pauli matrices divided by two, and let
\[
 P_* =\prod_{x\in\Lambda_L}\frac1{\ell!}
                  \sum_{\pi\in\mathfrak S_\ell}\mathcal U_{\pi,x}
\]
be the projection onto the tensor product of the symmetric site
subspaces.  Here $\mathcal U_{\pi,x}$ permutes the slots at $x$.
On this subspace $S_x^a=\sum_{i:\,x_i=x}s_i^a$ is the spin-$S$
representation.  In particular, if compression is needed for an individual
slot operator, then
\begin{equation}\label{eq:compressed-slot}
        P_*s_i^aP_* = \ell^{-1}S_{x_i}^aP_*.
\end{equation}

For each site bond use all $\ell^2$ slot edges above it.  On these
edges put independent Poisson exchanges of rate $1/2$ during
$[0,\beta)$, and at the seam from $\beta$ to $0$ insert an independent
uniform slot permutation at every site.  Let $\pi$ be the resulting
permutation of $\mathcal I$ and $\mathbb P_{\rm b}$ its underlying
picture law.  If $\mathcal E$ is the set of slot edges, then
\begin{equation}\label{eq:exchange-operator}
 e^{H_0}P_*=e^{\beta|\mathcal E|/4}
                    \E_{\rm b}\mathcal U_\pi,
 \qquad H_0=-\beta H_{S,L}.
\end{equation}
On the left $H_0$ is lifted to the slot space before compression.
Indeed,
$\boldsymbol s_i\cdot\boldsymbol s_j=\frac12\mathcal U_{(ij)}-\frac14I$,
and the mean chronological Poisson product has generator
$\frac12\sum_{e\in\mathcal E}(\mathcal U_e-I)$.
The all-slot sum above each bond commutes with $P_*$, while the seam
average inserts $P_*$.  This proves \eqref{eq:exchange-operator}.
This is T\'oth's exchange representation \cite{Toth1993}, with the
higher-spin construction of Nachtergaele
\cite[Section~2.2]{Nachtergaele1994}; see also
\cite[Proposition~2.1]{BL}.

If $A$ is a polynomial in the physical site spins, let $\widehat A$
be its slot lift obtained by replacing $S_x^a$ by
$\sum_{i:\,x_i=x}s_i^a$.  The physical trace is then
\begin{equation}\label{eq:physical-slot-trace}
       \Tr_{\rm phys}(Ae^{H_0})
          =\Tr_{\rm slot}(\widehat A e^{H_0}P_*).
\end{equation}
Following lines through repetitions of one sampled picture gives its
cycles.  In this slot trace with diagonal spin insertions a cycle has a constant
up or down color.  Without insertions its contribution is two.  We take
slot-time points immediately after an instantaneous operation, including
the seam.  In finite volume there are almost surely finitely many marks,
with distinct times avoiding any prescribed finite set of cuts.

We call a permutation law on a finite set \emph{admissible} if it is a
distributional limit of finite products of independent Bernoulli
transpositions and deterministic permutations.  Products here are ordered;
no commutation is assumed.  This class includes finite-time Poisson
exchanges at arbitrary nonnegative rates, uniform permutations on
specified subsets, and independent compositions of these operations.
For the Poisson assertion, divide time into small intervals and apply
each edge's Bernoulli indicator of at least one mark in a fixed order.
The approximation agrees with the picture unless an interval contains
at least two marks, whose probability tends to zero.  For a uniform
permutation, run positive-rate exchanges on every pair of the given
subset and let time tend to infinity.  This irreducible random walk on
the finite permutation group converges to its uniform invariant law.

We use the following finite stability facts.  A nonzero polynomial is
\emph{stable} if it does not vanish when all variables have positive
imaginary parts.  For every admissible law of a permutation $\pi$ of a
finite set $I$,
\begin{equation}\label{eq:stable-permutation-symbol}
              \E\prod_{i\in I}(u_i+v_{\pi(i)})
                    \quad\hbox{is stable}.
\end{equation}
Identification of variables preserves stability; differentiation, real
specialization, and bounded-degree coefficientwise limits preserve it
up to the zero polynomial.  Finally, if a random subset $A\subseteq I$
has a stable generating polynomial
$G_A(\boldsymbol z)=\E\prod_{j\in A}z_j$, then for every $J\subseteq I$,
\begin{equation}\label{eq:stable-presence}
            \Pr(J\subseteq A)\le\prod_{j\in J}\Pr(j\in A).
\end{equation}
These are the partial-symmetrization and negative-dependence facts in
\cite[Theorems~4.20 and~4.9]{BBL2009}; the exact finite formulations
and their elementary proofs appear in \cite[Lemmas~3.2--3.6]{Ord}.
In particular, the facts apply on a disjoint union of time layers,
without any spatial assumptions.

\subsection{Pins and exposure}

The finite constructions in this subsection allow independent Poisson
exchanges on an arbitrary finite slot graph, together with finitely many
prescribed deterministic permutations, Bernoulli transpositions, and
uniform permutations on prescribed subsets at fixed times.  All random operations are
independent under the base law.  Their order at
coincident prescribed times is fixed, and a slot-time point at such a
time lies after the entire ordered bundle of operations.  Thus extra stationary lines and
extra instantaneous exchanges are included.

A \emph{pin} is a specified slot-time point required to have up color.
For a finite deterministic pin set $E$, let $n_E$ be the number of
cycles avoiding $E$, let $\mathcal B_E$ be their union, and define
\begin{equation}\label{eq:pin-law}
 \mathcal Z_E=\E_{\rm b}2^{n_E},\qquad
       \mathbb P_E(d\omega)
         =\mathcal Z_E^{-1}2^{n_E(\omega)}\mathbb P_{\rm b}(d\omega).
\end{equation}
Cycles meeting pins are up; the others are colored independently and
fairly.  For a nonempty single-time set $T$ disjoint from $E$, put
$P=E\cup T$ and let $\sigma$ take each pin in $P$ to the next pin
on its directed cycle, at strictly positive elapsed time.  Set
\begin{equation}\label{eq:pin-return}
       r(E,T)=\frac1{|T|}\sum_{i\in T}
                              \mathbb P_P\{\sigma(i)\in T\}.
\end{equation}
The same sampled picture is repeated in defining $\sigma$.

To describe exposure, label lines at time zero and reveal all labels on
cycles missing $P$, their trajectories through a period, and all their
endpoint images.  The trajectories record their positions before and
after each elementary operation, including intermediate stages of an
ordered bundle at one time.  Denote these data by $\mathbf a$.  The slots not
occupied by exposed labels are called \emph{holes}.  For fixed exposure,
define a fresh law $\mathbb Q_{\mathbf a}$ on the holes as follows.
Exchanges whose two endpoints are holes have their original independent
Poisson clocks.  At a transposition exchanging an exposed label with a
hole, that hole moves into the vacated slot.  These forced transports are deterministic
bijections between successive hole layers.  A partly revealed uniform
permutation is completed by a uniform bijection between the remaining
source and target slots; a prescribed deterministic permutation restricts
to a deterministic bijection between those hole layers.
An elementary Bernoulli transposition touching
an exposed label has its outcome revealed; one touching only holes
retains its independent Bernoulli law.  Let $G_P$ denote the event that
every cycle of this fresh picture meets $P$.

\begin{lemma}[Finite pin inputs]\label{lem:pin-inputs}
For the finite picture above the following assertions hold.
\begin{enumerate}[label=(\roman*),leftmargin=*]
\item For every finite $E$ and nonempty single-time $T$ disjoint from $E$,
\begin{equation}\label{eq:simultaneous-pin}
                  \mathbb P_E(T\subseteq\mathcal B_E)
                                  \le(2r(E,T))^{|T|}.
\end{equation}
\item Almost surely under $\mathbb P_P$,
$g_P(\mathbf a):=\mathbb Q_{\mathbf a}(G_P)>0$, and the conditional
law of the remaining picture is
$\mathbb Q_{\mathbf a}(\,\cdot\mid G_P)$.
\item Let $\omega_{\rm m}$ be the Poisson mark counting measure and
let $\lambda_{\rm b}$ be edge rate times Lebesgue time measure.
For every integer $j\ge1$
and nonnegative measurable function $F$ of $j$ mark locations,
\begin{equation}\label{eq:factorial-domination}
 \E_E\sum_{(u_1,\ldots,u_j)\in\omega_{\rm m,\ne}^j}
                F(u_1,\ldots,u_j)
       \le 2^j\int F\,d\lambda_{\rm b}^{\otimes j}.
\end{equation}
The sum is over ordered distinct marks, and may be restricted to any
prescribed time order.  Independent additional fresh Poisson clocks
may be included with their own, undoubled, factorial intensity.
\end{enumerate}
\end{lemma}

\begin{proof}
The statements are the finite-pin forms of
\cite[Proposition~2.4, Proposition~2.5, and Lemma~2.3]{BL}.
We give the arguments to specify the scope of the extra operations and
the measure used after exposure.

Cut at all pin times and at zero, and let $\mathcal V$ be the disjoint
union of the slot layers at these cuts.  The maps between successive
cuts are independent admissible bijections.  Their symbols in disjoint
input and output variables multiply to give a stable polynomial
$\E_{\rm b}\prod_{i\in\mathcal V}(x_i+y_{\pi(i)})$.
Delete a nonpin $q$ by setting $x_q=y_q=t$ and extracting the
coefficient of $t$.  If its predecessor and successor are $i,j$, then
$(x_i+t)(t+y_j)$ is replaced by $x_i+y_j$.  A fixed point instead
contributes $2t$, hence a factor two.  Successive deletion therefore
gives the nonzero stable polynomial
\[
       \E_{\rm b}\left[2^{n_P}
                         \prod_{i\in P}(x_i+y_{\sigma(i)})\right].
\]
Differentiate in the input variables indexed by $P\setminus T$, set
all input variables to zero, and divide by $\mathcal Z_P$.
The result is
$\E_P\prod_{i\in T}y_{\sigma(i)}$, the stable generating polynomial
of the subset $\sigma(T)$.  By \eqref{eq:stable-presence} and the arithmetic--geometric
mean inequality,
\[
 \mathbb P_P\{\sigma(T)=T\}
  \le\prod_{j\in T}\mathbb P_P\{j\in\sigma(T)\}
  \le r(E,T)^{|T|}.
\]
Here $\sum_{j\in T}\mathbb P_P\{j\in\sigma(T)\}=|T|r(E,T)$
by bijectivity.  The event $\sigma(T)=T$ is precisely
$T\subseteq\mathcal B_E$: iterating the next-pin permutation from
$T$ then meets no pin in $E$.  Finally,
$2^{n_P}\le2^{n_E}\le2^{|T|}2^{n_P}$ and
$\mathcal Z_E\ge\mathcal Z_P$, proving (i).

For (ii), first fix a proposed set $A$ of initial exposed labels.
Use two independent Poisson clock families on every edge, each with
the original rate: accept the first family's clock when an endpoint
contains an $A$-label and the second when neither does.  This realizes
the base process.  Given the past, the superposed clock has twice the
total original rate, each edge is accepted with its original relative
rate, and the rejection probability is $1/2$.  The $A$-label paths
use only the first family.  Given these paths, the second family
restricted to the hole intervals remains fresh.  Revealing the
$A$-label images at a uniform permutation leaves a uniform residual
bijection.  At a Bernoulli transposition its outcome is revealed
exactly when an exposed label is involved; other outcomes remain
independent.  This constructs $\mathbb Q_{\mathbf a}$.

More explicitly, let $\mu_A$ be the base law of the paths and endpoint
images of the proposed labels.  Let $F_A$ mean that those paths avoid
$P$ and close on $A$, and let $c_A$ count their cycles on this event.
For every bounded measurable function $\Phi$ of the exposure and
remaining picture,
\begin{equation}\label{eq:exposure-branches}
 \E_P\Phi
   =\mathcal Z_P^{-1}\sum_A\int_{F_A}
       2^{c_A(\mathbf a)}
       \E_{\mathbb Q_{\mathbf a}}[\1_{G_P}\Phi]\,
                                                  \mu_A(d\mathbf a).
\end{equation}
Each full picture occurs in exactly one branch, namely its set of
labels on cycles missing $P$.  The exposure marginal on that branch
is therefore
$\mathcal Z_P^{-1}2^{c_A}g_P\,\mu_A$ restricted to $F_A$.
It assigns zero mass to $g_P=0$, and division by $g_P$ proves (ii).
In particular the exposure marginal includes the all-hit probability
$g_P$; it is not the unweighted law of the proposed histories.

For (iii), inserting one transposition joins two cycles or splits
one.  It increases the number of cycles missing the fixed pins by
at most one, and so multiplies the weight by at most two, regardless
of the other instantaneous operations.  The Poisson insertion
formula for $j$ ordered distinct marks, followed by division by
$\mathcal Z_E$, gives \eqref{eq:factorial-domination}.  Independent
fresh clocks have their ordinary factorial measure, proving the last
assertion.
\end{proof}

\subsection{The conditional next-pin first moment}

The conditioning in Lemma~\ref{lem:pin-inputs}(ii) couples entire
cycles.  The next result explains exactly why its next-pin first
moment can nevertheless be computed by a walk with fresh transitions.
Matrices in this subsection use the column-forward convention:
$\Pi e_i=e_{\pi(i)}$.

\begin{lemma}[Conditional next-pin kernel]\label{lem:next-pin}
Let $\Pi$ have an admissible permutation law on a finite set
$\mathcal V=I\sqcup D$, and put $K=\E\Pi$.  Let $G_I$ be the
event that every permutation cycle meets $I$.  On $G_I$, let
$\Sigma$ be the permutation matrix on $I$ taking each pin to its
next positive visit to $I$.  Then
\begin{equation}\label{eq:all-hit}
                    \Pr(G_I)=\det(I_D-K_{DD}).
\end{equation}
For every complex matrix $W$ on $I$, there is the polynomial identity
\begin{equation}\label{eq:next-pin-determinant}
 \E[\1_{G_I}\det(I_I-W\Sigma)]
  =\det\begin{pmatrix}
       I_I-WK_{II}&-WK_{ID}\\
       -K_{DI}&I_D-K_{DD}
       \end{pmatrix}.
\end{equation}
If $\Pr(G_I)>0$ and $I\ne\varnothing$, then $K_{DD}$ is transient
and
\begin{equation}\label{eq:next-pin-mean}
 \E[\Sigma\mid G_I]
    =K_{II}+K_{ID}(I_D-K_{DD})^{-1}K_{DI}.
\end{equation}
This is the first positive hitting kernel of $I$ for the Markov chain
with independent transitions $K$.  Empty determinants are one and
empty correction terms are zero.  In particular
\eqref{eq:next-pin-determinant}, and its linear coefficients in $W$,
remain valid when the all-hit probability is zero.
\end{lemma}

\begin{proof}
We reproduce the finite argument in \cite[Lemma~2.6 and
Proposition~2.7]{BL}, equivalently \cite[Lemma~4.2 and
Proposition~4.3]{Ord}.  Write $\bigwedge^k M$ for the action of $M$
on alternating $k$-tensors; its entries are the $k$-rowed minors of $M$.
For $0\le k\le|\mathcal V|$, exterior powers satisfy
\begin{equation}\label{eq:exterior-moments}
                \E\bigwedge\nolimits^k\Pi
                          =\bigwedge\nolimits^k K.
\end{equation}
Indeed, for a transposition $T_{ij}$,
$T_{ij}-I=-(e_i-e_j)(e_i-e_j)^*$ has rank one.  Every minor of
$I+u(T_{ij}-I)$ is affine in $u$, so its exterior power is the
mean exterior power of a Bernoulli transposition of probability $u$.
Independence and multiplicativity of exterior powers extend this
identity to products; continuity of minors extends it to their limits.

For any fixed matrix $B$, the principal-minor expansion now gives
\begin{equation}\label{eq:det-mean}
 \E\det(I-B\Pi)
  =\sum_k(-1)^k\Tr\left[(\bigwedge\nolimits^k B)
                            \E\bigwedge\nolimits^k\Pi\right]
  =\det(I-BK).
\end{equation}
Take $B=\diag(W,I_D)$.  If a permutation cycle is contained in
$D$, its indicator is a fixed vector of $B\Pi$, making the determinant
zero.  Otherwise $\Pi_{DD}$ is nilpotent, and elimination of this
block follows the successive nonpins on each cycle.  It yields
$\det(I-B\Pi)=\det(I_I-W\Sigma)$.
This proves \eqref{eq:next-pin-determinant}; setting $W=0$ proves
\eqref{eq:all-hit} without any invertibility assumption.

When the determinant is positive, take its Schur complement in
\eqref{eq:next-pin-determinant}.  Replace $W$ by $tW$ and compare
linear coefficients, using
$\det(I-tWA)=1-t\Tr(WA)+O(t^2)$.  Since $W$ is arbitrary,
the resulting trace identities recover every entry in
\eqref{eq:next-pin-mean}.
The nonnegative matrix $K_{DD}$ is substochastic.  Its spectral radius
is at most one; if it were one, the Perron eigenvalue would make
$I_D-K_{DD}$ singular.  Consequently
\[
                  (I_D-K_{DD})^{-1}=\sum_{n\ge0}K_{DD}^n.
\]
The term $K_{II}$ is a direct hit, and
$K_{ID}K_{DD}^nK_{DI}$ is a hit after $n+1$ intervening states in
$D$.  This proves the Markov interpretation.
\end{proof}

Apply the lemma at a fixed exposure by cutting at all pin times and
at the finitely many prescribed changes in hole sets.  Identify the
hole layers, which have equal cardinality.  The unconditioned interval
maps are independent admissible bijections; their direct sum followed
by the cyclic shift of the layers is an admissible permutation of the
layered set.  Its mean matrix gives the hole walk: jumps at the original
edge rates ($1/2$ for physical edges), prescribed deterministic restrictions
and forced transports, fresh residual uniform choices at each prescribed
uniform operation (including the seam), and fresh hole-only Bernoulli
outcomes.  Repeat this environment periodically, making independent
transition choices on every passage, including later passages through
the same physical time interval.  If $\tau_T^+$ is the first positive
visit to $T$ and $\tau_E$ the first visit to $E$, then
\begin{equation}\label{eq:fresh-return-identity}
 |T|r(E,T)=\E_P\sum_{i\in T}
              \Pr^{\rm walk}_{\mathbf a,i}(\tau_T^+<\tau_E).
\end{equation}
This follows by Lemma~\ref{lem:next-pin} under
$\mathbb Q_{\mathbf a}(\,\cdot\mid G_P)$ and then
Lemma~\ref{lem:pin-inputs}(ii).  The identity concerns a conditional
first moment.  The original repeated permutation paths retain their
cycle dependence; it is the determinant identity that replaces their
next-pin mean by fresh Markov transitions.

\subsection{Disk contractions and trace zero-freeness}

We next use \eqref{eq:stable-permutation-symbol} for complex matrix
insertions.  In any common qubit basis the polynomial of a permutation
layer, evaluated between input kets $(1,x_i)^T$ and algebraic output
covectors $(1,y_i)$, without complex conjugation, is
\begin{equation}\label{eq:disk-symbol}
       \E\prod_i(1+x_i y_{\pi(i)}).
\end{equation}
It has no zeros when all $|x_i|,|y_i|<1$.  To see this, substitute
\[
       u_i=i\frac{1+x_i}{1-x_i},\qquad
       v_j=i\frac{1-y_j}{1+y_j}
\]
in \eqref{eq:stable-permutation-symbol}.  Both families are in the
upper half-plane, and
$u_i+v_{\pi(i)}=2i(1+x_i y_{\pi(i)})/
[(1-x_i)(1+y_{\pi(i)})]$; the product of denominators is independent
of $\pi$.

\begin{lemma}[Zero-free layered traces]\label{lem:trace-zero-free}
On finitely many qubit slots, take a finite cyclic product of
permutation layers and tensor products of single-slot matrices,
with at least one permutation layer.
Each permutation layer is a nonzero scalar times the mean permutation
operator of an admissible law.  Between every two successive
permutation layers, suppose that the single-slot matrix on every slot
induces a projective map
\[
         (1,x)^T\longmapsto a(x)(1,\phi(x))^T
\]
with $a(x)\ne0$ and $|\phi(x)|<1$ for all $|x|\le1$.
Then the trace of the cyclic product is nonzero.
\end{lemma}

\begin{proof}
Ignore the nonzero scalar factors.  Group each single-slot layer with
the following permutation layer.  By compactness, its slot maps send
$|x|<1+\varepsilon$ into $|\phi(x)|<1$, with nonvanishing $a(x)$,
for some common $\varepsilon>0$.  Thus its multiaffine matrix
polynomial is nonzero for $|x_i|<1+\varepsilon$ and $|y_i|<1$,
by \eqref{eq:disk-symbol}.  Multiply the layer polynomials in disjoint
variables.

The trace is obtained by contracting each output variable with the
input variable on the matching slot of the next layer.  Such a
contraction acts on a bivariate section by
\begin{equation}\label{eq:trace-contraction}
                  A+Bx+Cy+Dxy\longmapsto A+D.
\end{equation}
To verify preservation of nonvanishing, fix all other variables in
their prescribed disks.  Choose $r_x<1+\varepsilon$ and $r_y<1$
with $r_xr_y>1$, and set $x=r_xz$, $y=r_yz$.  The resulting
polynomial has no zero in $|z|<1$ and has nonzero constant term $A$.
If its quadratic coefficient is nonzero, the product of its two
roots has modulus at least one, giving
$|D|r_xr_y\le|A|$.  If it has degree at most one, $D=0$.
In either case $|D|<|A|$ whenever $D\ne0$, so $A+D\ne0$.
After each contraction the polynomial stays multiaffine and nonzero
on the product of the remaining disks.  Contracting all pairs proves
the claim.  This is a disk form of the contraction method of
Asano and Ruelle \cite{Asano1970,Ruelle1971}, applied here to the
complete cyclic trace.
\end{proof}

Choose now the $s^2$ eigenbasis with vectors
$(1,\pm i)^T/\sqrt2$ in the $s^3$ coordinates.  The following
three slot maps will cover both the field comparison and the later
pressure argument.

\begin{lemma}[Single-slot maps]\label{lem:strip-maps}
In this basis, matrices in $SL_2(\R)$, real in the $s^3$ coordinates,
act by automorphisms of the unit disk.  If $\Rea z>0$, the matrix
$e^{zs^2}$ sends the closed disk strictly into the open disk.
There are absolute constants $b_*,a_0>0$ such that, whenever
$0<b\le b_*$, $f\in[-1,1]$, and $|\Ima z|<a_0$, the matrix
\begin{equation}\label{eq:strip-slot-map}
           e^{bs^2/2}\,e^{bzfs^3}\,e^{bs^2/2}
\end{equation}
also sends the closed disk strictly into the open disk, with its
projective coordinate defined throughout the closed disk.
\end{lemma}

\begin{proof}
The projective coordinate in the $s^3$ basis is
$\zeta=i(1-x)/(1+x)$.  It identifies the unit disk with the upper
half-plane.  A real determinant-one matrix acts there by an
upper-half-plane automorphism, proving the first assertion.
In the $s^2$ basis $e^{zs^2}$ acts by $x\mapsto e^{-z}x$,
proving the second.

For the last assertion, the rightmost factor first gives
$|x|\le r=e^{-b/2}$.  Its upper-half-plane image satisfies
\[
 \sin(\arg\zeta)=\frac{\Ima\zeta}{|\zeta|}
       =\frac{1-|x|^2}{|1-x|\,|1+x|}
       \ge\frac{1-r^2}{1+r^2}=\tanh(b/2).
\]
For $0<b\le b_*$, this is at least $cb$ with an absolute $c>0$.
Thus $\arg\zeta$ stays at angular distance at least $cb$ from
both $0$ and $\pi$.  The middle factor sends
$\zeta$ to $e^{-bzf}\zeta$, changing its argument by at most
$b|\Ima z|$.  Choose $a_0<c$; the image is still in the upper
half-plane.  The final factor shrinks the resulting disk coordinate
by $e^{-b/2}$.  None of these maps has a pole on the compact set
being mapped, proving the stated strictness and definedness.
\end{proof}

The next section uses a small positive field to supply pins.  The finite
identities proved here reduce uniform bounds on their avoidance to
estimates for the corresponding hole walk.
